% ============================================================
% COVER PAGE — an optional, full-bleed image page rendered before
% the title page. Not included by default; see README's "Cover
% page" section for how to turn this on.
%
% "Full-bleed" here means the image fills the entire physical page,
% edge to edge, with none of the document's normal 2cm margin —
% different from every other page in the book, which respects that
% margin. \newgeometry/\restoregeometry (both part of the geometry
% package, already loaded for the rest of this document) make that
% possible: margins drop to zero for this one page only, then get
% put back immediately after for everything that follows.
%
% This is "full-bleed within the page's own trim size" — the image
% fills the A4 page exactly, not a design that extends past A4's
% edges the way true press-ready bleed for a physically printed,
% trimmed book cover would. That distinction matters if this PDF is
% ever going to a professional print shop: press-ready bleed usually
% needs the cover treated as a separate file, built to that
% printer's own template (which specifies the exact bleed amount,
% commonly 3mm), not generated by this same LaTeX pipeline. For a
% PDF that's viewed on screen, printed at home/office, or sent to a
% simpler print-on-demand service that accepts a plain full-page
% image, what's below is the right tool.
%
% The image is stretched to fill \paperwidth x \paperheight exactly,
% not cropped — so it needs to already be sized to the page's own
% proportions (A4 portrait: roughly 0.707:1, width:height) to avoid
% visible distortion. cover-placeholder.png is sized correctly
% (1240x1754px) as a working example of that ratio.
% ============================================================
\newgeometry{margin=0pt}
\thispagestyle{empty}
\noindent\includegraphics[width=\paperwidth,height=\paperheight]{\giCoverPageImage}
\clearpage
\restoregeometry
